\section*{Descripción del problema de diseño}

En la gestión ganadera extensiva, el control del perímetro mediante cercados físicos tradicionales representa un elevado costo de instalación, mantenimiento constante e inflexibilidad para implementar sistemas de pastoreo rotacional eficiente. El problema de diseño consiste en definir y documentar la arquitectura de grado industrial de un \textbf{Collar Inteligente de Cerca Virtual y Monitoreo en Tiempo Real (DR-VF-01)}, capaz de delimitar potreros dinámicos mediante geolocalización satelital multiconstelación y emitir estímulos disuasivos progresivos cuando el bovino intenta sobrepasar los límites permitidos.

Para conciliar la entrega académica del semestre actual con la hoja de ruta del trabajo de tesis, el desarrollo se estructura formalmente en dos etapas:
\begin{itemize}[leftmargin=*]
    \item \textbf{Fase 1 (Entregable Semestre Actual -- Rev. 1.0):} Construcción de una tarjeta madre (\emph{Carrier Board}) con sockets modulares, centrada en el microcontrolador ESP32-S3, posicionamiento GPS u-blox MAX-M10S, telemetría LoRa SX1262 (915 MHz), gestión de recarga solar MPPT CN3791 (6V), medidor I$^2$C MAX17048 y emulación segura de la etapa disuasiva mediante optoacoplador VO617A con LED/Buzzer.
    \item \textbf{Fase 2 (Proyección de Tesis / Escalamiento Comercial):} Integración SMD a nivel de chip, conectividad híbrida con el módem celular Quectel BG95-M3 (LTE Cat-M1/NB-IoT), etapa elevadora de alto voltaje real ($2,5~\text{kV} / 0,2~\text{J}$ bajo norma IEC 60335-2-76) y encapsulado IP68.
\end{itemize}

\subsection*{Aspectos Críticos de Diseño}
\begin{itemize}[leftmargin=*]
    \item \textbf{Autonomía y Balance Energético MPPT:} Recarga fotovoltaica conmutada mediante el chip CN3791 regulado a $6~\text{V}$ y batería LiPo de $3,7~\text{V}$ ($1200\text{--}1500~\text{mAh}$ en Fase 1; ampliable a $3200~\text{mAh}$ en Fase 2).
    \item \textbf{Procesamiento en el Borde e IA (Edge AI):} Ejecución local en el microcontrolador ESP32-S3-N16R8 del algoritmo geométrico \emph{Point-in-Polygon} (a $10~\text{Hz}$) y clasificación de la conducta del animal mediante el núcleo de Machine Learning (MLC) de la IMU ST LSM6DSOX.
    \item \textbf{Seguridad y Aislamiento Galvánico:} Aislamiento de alto voltaje ($>3000~\text{V}_{\text{RMS}}$) entre la lógica de control ($3,3~\text{V}$) y el circuito disuasivo mediante el optoacoplador VO617A.
    \item \textbf{Costo Target de Grado Comercial:} Optimización de la lista de materiales (\emph{Target BOM}) por debajo de $\$50.00~\text{USD}$ ($\sim\$200.000~\text{COP}$) por unidad.
\end{itemize}

\section*{Descripción técnica del sistema}
La arquitectura funcional del collar se organiza en los siguientes bloques principales:
\begin{enumerate}[leftmargin=*]
    \item \textbf{Etapa de Cosecha Solar y Gestión de Potencia} \emph{[Fase 1]}: Panel solar monocristalino ($6~\text{V} / 1,5\text{--}3~\text{W}$), regulador MPPT Buck CN3791 (6V), medidor I$^2$C (MAX17048) y batería LiPo de $3,7~\text{V}$.
    \item \textbf{Unidad de Procesamiento Central (MCU)} \emph{[Fase 1]}: Microcontrolador ESP32-S3 Dual-Core (16 MB Flash, 8 MB PSRAM) a cargo de las geocercas, servidor web local (LittleFS) y punto de acceso Wi-Fi/BLE para configuración offline.
    \item \textbf{Subsistema de Posicionamiento Satelital GNSS} \emph{[Fase 1]}: Receptor multiconstelación u-blox MAX-M10S (GPS, GLONASS, Galileo, BeiDou) vía UART.
    \item \textbf{Subsistema de Telemetría Dual (LoRa + Celular IoT)}:
    \begin{itemize}
        \item \textbf{Semtech SX1262 (LoRa 915 MHz, $+22~\text{dBm}$) \emph{[Fase 1]}:} Transmisión primaria para zonas rurales.
        \item \textbf{Quectel BG95-M3 (LTE Cat-M1 / NB-IoT) \emph{[Expansión Futura -- Fase 2]}:} Módem celular de respaldo para reporte directo a la nube.
    \end{itemize}
    \item \textbf{Etapa de Estimulación Disuasiva Aislada}:
    \begin{itemize}
        \item \textbf{Emulación de Seguridad \emph{[Fase 1]}:} Optoacoplador VO617A con indicación por LED y Zumbador Piezoeléctrico.
        \item \textbf{Circuito Elevador de HV Real ($2,5~\text{kV} / 0,2~\text{J}$) \emph{[Fase 2]}:} Descarga controlada bajo norma IEC 60335-2-76.
    \end{itemize}
    \item \textbf{Subsistema Inercial e Inteligencia de Borde (IMU + MLC)} \emph{[Fase 1]}: ST LSM6DSOX para detección de movimiento (\emph{Wake-on-Motion}) y monitoreo de patrones de conducta.
\end{enumerate}